% GENERATED FILE. Edit canonical lessons, then run npm run generate:books.
% canonical-source-hash: fnv1a64:f1ecea255d8a5c7e
% canonical-lessons: JA-C397-minshuku, JA-C397-shukuhaku, JA-C397-ippaku, JA-C397-higaeri, JA-C397-azukaru

\chapter{A Room for the Night}
\label{ch:ja-a-room-for-the-night}
\hlchaptermodality{397}

\begin{chapteropening}
\textbf{By the end of this chapter:} \emph{I can book a room and a night's stay.}

Complete the last lesson of chapter 397: I can book a room and a night's stay.
\end{chapteropening}

\section[\texorpdfstring{\ja{みんしゅく} (minshuku)}{みんしゅく (minshuku)}]{\ja{みんしゅく} (minshuku) — a family-run guesthouse}

\label{lesson:JA-C397-minshuku}



\begin{quote}
Before the new one: say the Japanese for strolling, then the Japanese for wandering restlessly.
\end{quote}

\subsection*{You'll want to know: \ja{みんしゅく}}
\textbf{\ja{みんしゅく}} — \emph{minshuku} — \textquotedblleft{}a family-run guesthouse\textquotedblright{}.

\begin{quote}
\textbf{\ja{うみの} \ja{ちかくに} \ja{みんしゅくが} \ja{あります。}} — \emph{umi no chikaku ni minshuku ga arimasu} — \textquotedblleft{}There is a family-run guesthouse near the sea.\textquotedblright{}
\end{quote}

\begin{cousinweb}
\textbf{\ja{みんしゅく}} (\emph{minshuku}) is Sino-Japanese (kango): two Chinese characters read with their on'yomi, readings borrowed from Chinese. They are \emph{min}, \textquotedblleft{}people\textquotedblright{} (Middle Chinese \emph{mjin}) and \emph{shuku}, \textquotedblleft{}inn\textquotedblright{} (\emph{sjuwk}) (KANJIDIC2; Middle Chinese in Baxter's transcription).
\end{cousinweb}

\begin{grammarlens}[title={What you've built}]
One more word for getting around.
\end{grammarlens}

\subsection*{Your turn}
\begin{itemize}
\raggedright
  \item \emph{Say it:} \emph{minshuku}
  \item \emph{Say it:} \emph{minshuku}, once more
  \item \emph{Say it:} \emph{umi no chikaku ni minshuku ga arimasu} — There is a family-run guesthouse near the sea.
  \item \emph{Recall:} say the Japanese for strolling, then the Japanese for wandering restlessly, then say \emph{minshuku} again
\end{itemize}

\begin{tcolorbox}[breakable,colback=teal!4,colframe=teal!35!black,title={Before you move on}]
What does \ja{みんしゅく} mean? (\textquotedblleft{}A family-run guesthouse\textquotedblright{}.) Say it once more.
\end{tcolorbox}
\section[\texorpdfstring{\ja{しゅくはく} (shukuhaku)}{しゅくはく (shukuhaku)}]{\ja{しゅくはく} (shukuhaku) — accommodation; lodging}

\label{lesson:JA-C397-shukuhaku}



\begin{quote}
Before the new one: say the Japanese for wandering restlessly, then the Japanese for a family-run guesthouse.
\end{quote}

\subsection*{You'll want to know: \ja{しゅくはく}}
\textbf{\ja{しゅくはく}} — \emph{shukuhaku} — \textquotedblleft{}accommodation; lodging\textquotedblright{}.

\begin{quote}
\textbf{\ja{しゅくはくの} \ja{よやくを} \ja{しました。}} — \emph{shukuhaku no yoyaku o shimashita} — \textquotedblleft{}I booked somewhere to stay.\textquotedblright{}
\end{quote}

\begin{cousinweb}
\textbf{\ja{しゅくはく}} (\emph{shukuhaku}) is Sino-Japanese (kango): two Chinese characters read with their on'yomi, readings borrowed from Chinese. They are \emph{shuku}, \textquotedblleft{}lodging\textquotedblright{} (Middle Chinese \emph{sjuwk}) and \emph{haku}, \textquotedblleft{}overnight stay\textquotedblright{} (\emph{bak}) (KANJIDIC2; Middle Chinese in Baxter's transcription).
\end{cousinweb}

\begin{grammarlens}[title={What you've built}]
One more word for getting around.
\end{grammarlens}

\subsection*{Your turn}
\begin{itemize}
\raggedright
  \item \emph{Say it:} \emph{shukuhaku}
  \item \emph{Say it:} \emph{shukuhaku}, once more
  \item \emph{Say it:} \emph{shukuhaku no yoyaku o shimashita} — I booked somewhere to stay.
  \item \emph{Recall:} say the Japanese for wandering restlessly, then the Japanese for a family-run guesthouse, then say \emph{shukuhaku} again
\end{itemize}

\begin{tcolorbox}[breakable,colback=teal!4,colframe=teal!35!black,title={Before you move on}]
What does \ja{しゅくはく} mean? (\textquotedblleft{}Accommodation; lodging\textquotedblright{}.) Say it once more.
\end{tcolorbox}
\section[\texorpdfstring{\ja{いっぱく} (ippaku)}{いっぱく (ippaku)}]{\ja{いっぱく} (ippaku) — a one-night stay}

\label{lesson:JA-C397-ippaku}



\begin{quote}
Before the new one: say the Japanese for a family-run guesthouse, then the Japanese for accommodation.
\end{quote}

\subsection*{You'll want to know: \ja{いっぱく}}
\textbf{\ja{いっぱく}} — \emph{ippaku} — \textquotedblleft{}a one-night stay\textquotedblright{}.

\begin{quote}
\textbf{\ja{りょかんに} \ja{いっぱくしました。}} — \emph{ryokan ni ippaku shimashita} — \textquotedblleft{}I stayed one night at a Japanese inn.\textquotedblright{}
\end{quote}

\begin{cousinweb}
\textbf{\ja{いっぱく}} (\emph{ippaku}) is Sino-Japanese (kango): two Chinese characters read with their on'yomi, readings borrowed from Chinese. They are \emph{ichi}, \textquotedblleft{}one\textquotedblright{} (Middle Chinese \emph{'jit}) and \emph{haku}, \textquotedblleft{}overnight stay\textquotedblright{} (\emph{bak}), joined as \emph{ippaku} (KANJIDIC2; Middle Chinese in Baxter's transcription).
\end{cousinweb}

\begin{grammarlens}[title={What you've built}]
One more word for getting around.
\end{grammarlens}

\subsection*{Your turn}
\begin{itemize}
\raggedright
  \item \emph{Say it:} \emph{ippaku}
  \item \emph{Say it:} \emph{ippaku}, once more
  \item \emph{Say it:} \emph{ryokan ni ippaku shimashita} — I stayed one night at a Japanese inn.
  \item \emph{Recall:} say the Japanese for a family-run guesthouse, then the Japanese for accommodation, then say \emph{ippaku} again
\end{itemize}

\begin{tcolorbox}[breakable,colback=teal!4,colframe=teal!35!black,title={Before you move on}]
What does \ja{いっぱく} mean? (\textquotedblleft{}A one-night stay\textquotedblright{}.) Say it once more.
\end{tcolorbox}
\section[\texorpdfstring{\ja{ひがえり} (higaeri)}{ひがえり (higaeri)}]{\ja{ひがえり} (higaeri) — a day trip}

\label{lesson:JA-C397-higaeri}



\begin{quote}
Before the new one: say the Japanese for accommodation, then the Japanese for a one-night stay.
\end{quote}

\subsection*{You'll want to know: \ja{ひがえり}}
\textbf{\ja{ひがえり}} — \emph{higaeri} — \textquotedblleft{}a day trip\textquotedblright{}.

\begin{quote}
\textbf{\ja{やまに} \ja{ひがえりで} \ja{いきました。}} — \emph{yama ni higaeri de ikimashita} — \textquotedblleft{}We went to the mountains on a day trip.\textquotedblright{}
\end{quote}

\begin{cousinweb}
\textbf{\ja{ひがえり}} (\emph{higaeri}) is a native Japanese word (wago), inherited rather than borrowed from Chinese, built from \emph{hi}, \textquotedblleft{}day\textquotedblright{} and \emph{gaeri}, \textquotedblleft{}returning\textquotedblright{} (from \emph{kaeru}); its kanji were attached later as kun'yomi (JMdict).
\end{cousinweb}

\begin{grammarlens}[title={What you've built}]
One more word for getting around.
\end{grammarlens}

\subsection*{Your turn}
\begin{itemize}
\raggedright
  \item \emph{Say it:} \emph{higaeri}
  \item \emph{Say it:} \emph{higaeri}, once more
  \item \emph{Say it:} \emph{yama ni higaeri de ikimashita} — We went to the mountains on a day trip.
  \item \emph{Recall:} say the Japanese for accommodation, then the Japanese for a one-night stay, then say \emph{higaeri} again
\end{itemize}

\begin{tcolorbox}[breakable,colback=teal!4,colframe=teal!35!black,title={Before you move on}]
What does \ja{ひがえり} mean? (\textquotedblleft{}A day trip\textquotedblright{}.) Say it once more.
\end{tcolorbox}
\section[\texorpdfstring{\ja{あずかる} (azukaru)}{あずかる (azukaru)}]{\ja{あずかる} (azukaru) — to look after (something for someone)}

\label{lesson:JA-C397-azukaru}



\begin{quote}
Before the new one: say the Japanese for a one-night stay, then the Japanese for a day trip.
\end{quote}

\subsection*{You'll want to know: \ja{あずかる}}
\textbf{\ja{あずかる}} — \emph{azukaru} — \textquotedblleft{}to look after (something for someone)\textquotedblright{}.

\begin{quote}
\textbf{\ja{ともだちの} \ja{にもつを} \ja{あずかりました。}} — \emph{tomodachi no nimotsu o azukarimashita} — \textquotedblleft{}I looked after my friend's luggage.\textquotedblright{}
\end{quote}

Polite \textbf{\ja{あずかります}}, \emph{azukarimasu}; a godan verb. Past (polite): \textbf{\ja{りょかんが} \ja{にもつを} \ja{あずかりました。}}, \emph{ryokan ga nimotsu o azukarimashita}.

\begin{cousinweb}
\textbf{\ja{あずかる}} (\emph{azukaru}) is a native Japanese word (wago), inherited rather than borrowed from Chinese; its kanji was attached later as a kun'yomi (JMdict). It is the partner of \emph{azukeru}, \textquotedblleft{}to deposit\textquotedblright{} (JMdict).
\end{cousinweb}

\begin{grammarlens}[title={What you've built}]
One more word for getting around.
\end{grammarlens}

\subsection*{Your turn}
\begin{itemize}
\raggedright
  \item \emph{Say it:} \emph{azukaru}
  \item \emph{Say it:} \emph{azukaru}, once more
  \item \emph{Say it:} \emph{tomodachi no nimotsu o azukarimashita} — I looked after my friend's luggage.
  \item \emph{Recall:} say the Japanese for a one-night stay, then the Japanese for a day trip, then say \emph{azukaru} again
\end{itemize}

\begin{tcolorbox}[breakable,colback=teal!4,colframe=teal!35!black,title={Before you move on}]
What does \ja{あずかる} mean? (\textquotedblleft{}To look after (something for someone)\textquotedblright{}.) Say it once more.
\end{tcolorbox}
